% Abhakseite „Das kann ich“ – Zone nur im Gesamt (\mitzone)
%% TODO Ich-kann-Sätze: Platzhalter wie in den Aufgabendateien
\begin{abhakseite}
\ifdefined\mitzone
\abhakgruppe{Kennst du schon}
\abhak{1}{Skalarprodukt berechnen und „null heißt senkrecht“ anwenden}
\abhak{2}{Richtungsvektoren bilden und Beträge berechnen}
\abhak{3}{Lineare und quadratische Gleichungen sowie kleine Gleichungssysteme lösen, Lösungen gegen einen Bereich sieben}
\abhak{4}{Kollinearität prüfen (Vielfaches)}
\abhak{5}{Normalenvektor aus der Koordinatengleichung ablesen, Spannvektoren aus der Parameterform lesen}
\abhak{6}{Ähnlichkeit rechtwinkliger Dreiecke und Seitenverhältnisse}
\abhak{7}{Richtungsvektoren bilden und Beträge berechnen – Fehler finden}
\abhak{8}{Richtungsvektoren bilden und Beträge berechnen – selbst rechnen}
\fi
\abhakgruppe{Einheit 1 · Rechte Winkel an Dreiecken nachweisen}
\abhak{9}{„Welches Paar ist senkrecht – und welche Prüfregel gilt?“}
\abhak{10}{„Null zeigen oder null setzen?“}
\abhak{11}{„Für alle oder für eins?“}
\abhak{12}{Rechte Winkel nachweisen}
\abhak{13}{Rechte Winkel nachweisen – Fehler finden, Begründen, Anwendung}
\abhakgruppe{Einheit 2 · Rechte Winkel rückwärts}
\abhak{14}{Rechte Winkel rückwärts}
\abhak{15}{Rechte Winkel rückwärts – Fehler finden, Begründen, Anwendung}
\abhakgruppe{Einheit 3 · Senkrecht zu Geraden und Ebenen}
\abhak{16}{Senkrecht zu Geraden und Ebenen}
\abhak{17}{Senkrecht zu Geraden und Ebenen – Fehler finden, Begründen, Anwendung}
\abhakgruppe{Einheit 4 · Lot und Extremum}
\abhak{18}{Lot und Extremum}
\abhak{19}{Flächengleichung des flächenkleinsten Dreiecks begründen}
\abhak{20}{Lot und Extremum – Fehler finden, Begründen, Anwendung}
\end{abhakseite}
